%==============================================================================
% ABSTRAK
%==============================================================================
\chapter*{ABSTRAK}
\addcontentsline{toc}{chapter}{ABSTRAK}

Penelitian mengenai topik ini dilaksanakan di P.T. Bangga Teknologi Produktif (Karla Bionics) yang merupakan hasil riset dan pengabdian masyarakat oleh akademisi Institut Teknologi Bandung (ITB), memposisikan diri sebagai penggerak inovasi sosial yang menekankan keberlanjutan dan dialog dengan para pihak berkepentingan, sehingga setiap solusi yang dihasilkan sesuai dengan kebutuhan masyarakat. Karla Bionics ingin berkontribusi pada transformasi digital layanan kesehatan dengan mendukung mitra rumah sakit penyelenggara Program Jaminan Kesehatan Nasional (JKN) BPJS Kesehatan, yang menghadapi permasalahan klaim sering berstatus \textit{pending} akibat ketidaktepatan pemberian kode diagnosis (ICD-10) dan prosedur (ICD-9-CM). Kondisi tersebut menunda pembayaran klaim, mengganggu arus kas rumah sakit, dan menambah beban kerja petugas koding, proses koding yang masih manual dan bergantung pada pengalaman petugas juga membuat hasilnya kurang konsisten. Untuk menjawab permasalahan tersebut, sistem pendukung keputusan berbasis web bernama INA-CBG \textit{Coding Assistant} dikembangkan untuk membantu petugas koding mengunggah resume medis, mengekstraksi teks dokumen melalui \textit{parsing} lokal dengan \textit{fallback} OCR berbasis kecerdasan buatan, memperoleh rekomendasi kode ICD-10 dan ICD-9-CM beserta bukti pendukung dan tingkat keyakinannya, mengambil keputusan terhadap kode tersebut (konfirmasi, tolak, atau tambah manual), serta menghitung estimasi grup dan tarif INA-CBG. Sistem dibangun dengan \textit{backend} FastAPI, \textit{worker} Celery melalui antrian Redis, basis data PostgreSQL, penyimpanan berkas Cloudflare R2, dan integrasi Google Gemini yang dilengkapi mekanisme \textit{auto-failover} ke OpenRouter. Sistem juga menyembunyikan identitas pasien sebelum data dikirim ke penyedia AI eksternal dan mencatat seluruh aktivitas penting pada log audit. Hasil pengujian \textit{blackbox} menunjukkan seluruh skenario uji berjalan dan lulus.

\vspace{12pt}
\noindent Kata kunci: \textit{pending claim}, koding medis, ICD-10, INA-CBG, sistem pendukung keputusan
